\section{Cavity inference on a declared factor graph}

Matrix scaling solved a local problem: given one pair score and two marginals, construct a
probability-exact pair law.  The preceding compatibility example showed that solving this problem
edge by edge does not guarantee one joint law for all vertices.  Cavity inference starts after that
additional joint factorization has been declared.  It propagates the effect of the rest of the
factor graph into a local marginal, thereby separating the pair factor, the inference message, and
the resulting pair law.

The cavity method begins from a joint factorization.  For a pairwise categorical model on a graph
$G=(V,E)$, write
\begin{equation}
P(x)=\frac1Z\prod_{i\in V}\psi_i(x_i)
\prod_{(i,j)\in E}\psi_{ij}(x_i,x_j).
\label{eq:book-pair-factor-law}
\end{equation}
On a tree, deleting the factor on $(i,j)$ separates the two sides.  The cavity message carried by
the directed edge $i\to j$ obeys
\begin{equation}
m_{i\to j}(a)
\propto
\psi_i(a)
\prod_{k\in\partial i\setminus\{j\}}
\sum_b\psi_{ik}(a,b)m_{k\to i}(b).
\label{eq:book-cavity-message}
\end{equation}
The exact pair marginal is then
\begin{equation}
P_{ij}(a,b)
\propto
\psi_{ij}(a,b)m_{i\to j}(a)m_{j\to i}(b).
\label{eq:book-cavity-pair-marginal}
\end{equation}
These equations are exact on trees and form the Bethe approximation on loopy graphs
\cite{mezard1987,mezardmontanari2009}.

One edge now carries several related objects:
\begin{center}
\small
\begin{tabular}{@{}>{\raggedright\arraybackslash}p{0.23\linewidth}
                >{\raggedright\arraybackslash}p{0.30\linewidth}
                >{\raggedright\arraybackslash}p{0.35\linewidth}@{}}
\toprule
Object & Definition & Meaning \\
\midrule
$\log\psi_{ij}(a,b)$ & factor parameter, projected into a gauge when needed & direct pair term in the declared joint law \\
$m_{i\to j}(a)$ & marginal state of $i$ with the $(i,j)$ factor removed & directed inference message \\
$P_{ij}(a,b)$ & compatible marginal after messages are recombined & exact or approximate pair law \\
$\chi_{i\leftarrow j}$ & derivative of a marginal with respect to a field & propagated statistical susceptibility \\
$G_{i\leftarrow j}(C)$ & finite endpoint substitution at fixed background & background-conditioned model response \\
\bottomrule
\end{tabular}
\end{center}

The distinction matters even when all five objects are computed from one Potts model.  The factor
is static.  Messages vary with boundary conditions and the rest of the graph.  Susceptibility sums
all paths permitted by the joint law.  The finite-background audit conditions on one realized
context rather than marginalizing it.  Agreement among these objects is a theorem only under the
corresponding factorization, topology, and conditioning operation.

There is also a terminology boundary.  Graphical-model cavity inference deletes an edge or factor
while retaining an endpoint variable as the argument of a message.  The endpoint-excluding audit
in this book withholds the realized identities of both $X_i$ and $X_j$ from support selection and
reference construction.  Both use deletion to control information flow, but they delete different
objects and license different conditional statements.  We reserve \emph{cavity message} for
Eq.~\ref{eq:book-cavity-message} and use \emph{endpoint exclusion} for the architectural
measurability condition.

This supplies a controlled protein benchmark.  Generate categorical sequences on a tree from a
known Potts law, then report the projected factor $G_{ij}^{\mathrm P}$, the BP marginal, the
susceptibility, and the finite substitution response in common anchor coordinates.  Exact
tree identities test the implementation.  Adding loops exposes the Bethe residual, which can be
measured by marginal error, calibration, or discrepancy from exact enumeration on small systems.
For dense protein-family models, tree exactness is unavailable unless a sparse factorization has
been declared and validated.

Replica-symmetry breaking moves the object one level higher: a single cavity message is replaced
by a distribution over messages or pure states.  Survey propagation uses this idea to derive
algorithms for random constraint problems \cite{mezard2002sat}.  In the language of this book it
suggests an ensemble of operator graphs or message fields, with uncertainty attached to the edge
object itself.  That extension belongs to the open problem of graph ensembles rather than to a
single deterministic interaction graph.
